\documentclass[11pt, a4paper]{article}

\usepackage[T1]{fontenc}
\usepackage[utf8]{inputenc}
\usepackage{tgpagella}
\usepackage[spanish, es-noshorthands]{babel}
\usepackage{amsmath, amssymb, bm}
\usepackage{tabularx}
\usepackage{booktabs}
\usepackage[most]{tcolorbox}
\usepackage[margin=2.2cm]{geometry}
\usepackage{ragged2e}
\usepackage{fancyhdr}

\pagestyle{fancy}
\fancyhf{}
\setlength{\headheight}{16pt}
\lhead{\textbf{UCB Sede Tarija} | Ing. Mecatrónica}
\chead{}
\rhead{IMT-121 Circuitos Electrónicos I | Plantilla Informe IEEE}
\lfoot{Prof: M.Sc. Ing. Bernardo Quiroga Turdera}
\rfoot{Página \thepage}
\renewcommand{\headrulewidth}{0.4pt}

\definecolor{ucbblue}{RGB}{0, 51, 102}

\begin{document}

\begin{tcolorbox}[colback=ucbblue!5!white, colframe=ucbblue, title=\textbf{Artefacto A3: Estructura Editorial Compatible IEEE (EC3)}]
    \centering \textbf{Caracterización Analítica, Computacional y Experimental de Circuitos RC y RLC}
\end{tcolorbox}

\noindent\textit{Aviso Formativo: Este documento es el esquema de contenidos para el informe final de EC3 (30\% de evaluación). Deberá vaciarse en la plantilla institucional de dos columnas. Fecha propuesta de entrega: 06/11/2026, 23:59 hrs.}

\section*{Estructura de Secciones Exigidas}

\textbf{Autores y Filiación:} Nombres reales, correos electrónicos. No inventar datos. \\
\textbf{Resumen (Abstract):} 120–200 palabras. Indique claramente qué mediciones lograron efectuarse en el banco. Si hubo un NO-GO experimental que bloqueó el circuito RLC, indíquelo como una limitación honesta de la investigación. \\
\textbf{Palabras Clave:} RC, RLC, respuesta transitoria, LTspice, constante de tiempo, amortiguamiento, osciloscopio.

\subsubsection*{I. Introducción}
Describa el objetivo de la validación y por qué es necesario confrontar una predicción en hoja, un simulador (LTspice) y el montaje del "Mundo Real". Declare el alcance efectivo de las pruebas.

\subsubsection*{II. Fundamento Teórico}
Plantee de forma concisa el modelo de inercia y su diferencial (RC: $\tau = R_{eq}C$). Para el circuito RLC describa el surgimiento de la ecuación de 2do Orden ($LC v_C'' + R_{\text{tot}}C v_C' + v_C = v_s$), defina rigurosamente $\omega_0$ y el factor adimensional $\zeta$. Citar bibliografía consultada.

\subsubsection*{III. Metodología}
Tabule los componentes reales empleados. Especifique el modelo comercial del Osciloscopio y Generador. Declare la frecuencia, $V_{pp}$, $V_{\text{offset}}$ y modo de disparo usados. Detalle la topología de la puesta a tierra de sus sondas para certificar que el método es seguro. Si la práctica sufrió un "NO-GO", debe relatarse qué lo detonó.

\subsubsection*{IV. Resultados y Tablas (Contraste Numérico)}
Enfoque en la trazabilidad. La Tabla de Resultados \textbf{debe comparar} el Modelo Teórico vs LTspice vs Experimental. Use "NO MEDIDO" si no obtuvo datos por NO-GO, pero jamás invente o copie resultados simulados en las columnas del osciloscopio.

\vspace{0.15cm}
\renewcommand{\arraystretch}{1.3}
\noindent
\begin{tabularx}{\textwidth}{
  >{\bfseries\RaggedRight\hsize=1.1\hsize}X 
  >{\RaggedRight\arraybackslash\hsize=1.0\hsize}X 
  >{\RaggedRight\arraybackslash\hsize=1.0\hsize}X 
  >{\RaggedRight\arraybackslash\hsize=1.1\hsize}X 
  >{\RaggedRight\arraybackslash\hsize=0.8\hsize}X}
\toprule
\textbf{Cantidad} & \textbf{Analítico} & \textbf{LTspice} & \textbf{Experimental (Si GO)} & \textbf{Unidad} \\
\midrule
RC $\tau$     & $10$ (nominal) & \textit{Simulado} & \textit{Medido o N/D} & ms \\
RLC $\zeta$   & $0.8$ (\textit{Ejemplo}) & \textit{No aplica} & \textit{Inferido o NO MED} & -- \\
\bottomrule
\end{tabularx}

\subsubsection*{V. Discusión de Discrepancias}
Realice la disección del error relativo ($\varepsilon$). Por ejemplo: Si su $\tau_{\text{med}} = 10.7\,\text{ms}$ y el teórico era $10.0\,\text{ms}$, discuta matemáticamente si la tolerancia del 5\% de su capacitor y los probables $50\,\Omega$ de resistencia de salida del generador $R_{\text{gen}}$ justifican y validan esta diferencia, en lugar de atribuirlo ciegamente a un "error de medición".

\subsubsection*{VI. Conclusiones y Anexos}
Presente de 3 a 5 aseveraciones categóricas basadas únicamente en sus tablas de datos. Incluya en Anexos las capturas de pantalla reales (USB del osciloscopio) identificando correctamente las trazas y escalas.

\vspace{0.3cm}
\begin{tcolorbox}[colback=gray!5, colframe=black, title=\textbf{Pauta de Evaluación Editorial Propuesta (30\% de EC3)}]
\small (1) Reproducibilidad y Criterio de Método (25\%) | (2) Datos y Trazabilidad Auténtica (25\%) | (3) Análisis Cuantitativo y de Error (25\%) | (4) Discusión y Justificación Causal (20\%) | (5) Presentación y Citas IEEE (5\%).
\end{tcolorbox}

\end{document}
